\documentclass[11pt]{article}
\usepackage[a4paper,margin=1in]{geometry}
\usepackage{amsmath,amssymb}
\usepackage{hyperref}
\title{Canonical Data Model for PyTex}
\author{PyTex Contributors}
\date{\today}

\begin{document}
\maketitle

\section{Motivation}

Texture and diffraction workflows are sensitive to convention drift. A vector without an attached frame, a plane without a declared reciprocal basis, or an orientation without explicit crystal and specimen frames is scientifically under-specified.

\section{Design}

PyTex therefore defines stable primitives for:

\begin{itemize}
  \item reference frames and transforms,
  \item symmetry specifications,
  \item lattice and basis representations,
  \item rotations, orientations, and misorientations,
  \item texture-domain containers,
  \item EBSD-domain containers,
  \item diffraction-domain containers,
  \item provenance records.
\end{itemize}

\section{Canonical Convention}

The initial canonical convention set is:

\begin{align*}
  \text{handedness} &= \text{right-handed}, \\
  q &= (w, x, y, z), \\
  \text{Euler labels} &= (\phi_1, \Phi, \phi_2), \\
  \mathbf{a}_i^\ast \cdot \mathbf{a}_j &= \delta_{ij}.
\end{align*}

\section{Implementation Rule}

Public APIs should not expose raw arrays where domain types are needed to recover scientific meaning.

\end{document}
